\section{Foundations of probability theory}\label{sec:probability}

We recall the measure-theoretic foundations of probability, following Kolmogorov's axiomatization \cite{Billingsley,Durrett}. The purpose of this section is twofold: to fix notation used throughout, and to exhibit the structures---$\sigma$-algebras, measures, conditional expectations---that the Fourier--Mellin duality will later transform.

\begin{definition}[Probability space]\label{def:prob-space}
A \emph{probability space} is a triple $(\Omega, \cF, P)$ consisting of:
\begin{enumerate}[label=(\roman*)]
\item a non-empty set $\Omega$ (the \emph{sample space}),
\item a $\sigma$-algebra $\cF \subseteq 2^\Omega$ (the \emph{event algebra}),
\item a measure $P : \cF \to [0,1]$ with $P(\Omega) = 1$ (the \emph{probability measure}).
\end{enumerate}
An element $\omega \in \Omega$ is an \emph{outcome}; a set $A \in \cF$ is an \emph{event}.
\end{definition}

\begin{remark}
The axioms encode three intuitive requirements. First, $\cF$ is closed under countable set operations, so any statement built from countably many events remains an event. Second, $P(\varnothing) = 0$ and $P(\Omega) = 1$, normalizing certainty. Third, $P$ is \emph{countably additive}: for pairwise disjoint events $A_1, A_2, \ldots \in \cF$,
\[
P\!\left(\bigcup_{n=1}^\infty A_n\right) = \sum_{n=1}^\infty P(A_n).
\]
Countable additivity is the key axiom; it is what makes the theory non-trivial, distinguishing probability from finitely additive ``charges'' and enabling limiting arguments such as the CLT.
\end{remark}

\begin{definition}[Random variable]\label{def:rv}
A \emph{random variable} on $(\Omega, \cF, P)$ is a measurable function $X : (\Omega, \cF) \to (\R, \mathcal{B}(\R))$, where $\mathcal{B}(\R)$ denotes the Borel $\sigma$-algebra. More generally, for a measurable space $(S, \mathcal{S})$, an \emph{$S$-valued random variable} is a measurable map $X : (\Omega, \cF) \to (S, \mathcal{S})$. The \emph{law} (or \emph{distribution}) of $X$ is the pushforward measure $\mu_X = P \circ X^{-1}$ on $(S, \mathcal{S})$.
\end{definition}

\begin{definition}[Conditional probability]\label{def:cond-prob}
Let $(\Omega, \cF, P)$ be a probability space and let $B \in \cF$ with $P(B) > 0$. The \emph{conditional probability} of $A \in \cF$ given $B$ is
\[
P(A \mid B) \;=\; \frac{P(A \cap B)}{P(B)}.
\]
\end{definition}

\begin{remark}
Conditioning on an event $B$ with $P(B) > 0$ restricts the probability space: the map $A \mapsto P(A \mid B)$ is itself a probability measure on $(\Omega, \cF)$, concentrated on $B$. The definition encodes a selection effect---knowing that $B$ has occurred reweights the sample space.
\end{remark}

\begin{proposition}[Law of total probability]\label{prop:total-prob}
Let $B_1, B_2, \ldots$ be a countable partition of $\Omega$ into events of positive probability. Then for any $A \in \cF$,
\[
P(A) = \sum_{n=1}^\infty P(A \mid B_n)\, P(B_n).
\]
\end{proposition}

\begin{proof}
Since $A = \bigcup_n (A \cap B_n)$ with the sets $A \cap B_n$ pairwise disjoint,
\[
P(A) = \sum_n P(A \cap B_n) = \sum_n P(A \mid B_n)\, P(B_n). \qedhere
\]
\end{proof}

\begin{theorem}[Bayes' rule]\label{thm:bayes}
Let $A, B \in \cF$ with $P(A) > 0$ and $P(B) > 0$. Then
\begin{equation}\label{eq:bayes}
P(A \mid B) \;=\; \frac{P(B \mid A)\, P(A)}{P(B)}.
\end{equation}
More generally, if $\{A_n\}$ is a countable partition of $\Omega$ with $P(A_n) > 0$ for each $n$, then
\begin{equation}\label{eq:bayes-partition}
P(A_k \mid B) \;=\; \frac{P(B \mid A_k)\, P(A_k)}{\displaystyle\sum_n P(B \mid A_n)\, P(A_n)}.
\end{equation}
\end{theorem}

\begin{proof}
Equation~\eqref{eq:bayes} follows from the symmetry $P(A \cap B) = P(B \mid A)\, P(A) = P(A \mid B)\, P(B)$. Substituting the law of total probability (\cref{prop:total-prob}) into the denominator yields \eqref{eq:bayes-partition}.
\end{proof}

\begin{remark}[Bayesian interpretation]
Bayes' rule inverts the direction of conditioning. In the Bayesian reading, $P(A_k)$ is the \emph{prior} probability of hypothesis $A_k$; $P(B \mid A_k)$ is the \emph{likelihood} of the observed data $B$ under hypothesis $A_k$; and $P(A_k \mid B)$ is the \emph{posterior} probability of $A_k$ given the data. The denominator $P(B) = \sum_n P(B \mid A_n)\, P(A_n)$ is a normalizing constant (the \emph{marginal likelihood} or \emph{evidence}). Bayes' rule is thus the mechanism by which data transforms prior beliefs into posterior beliefs.
\end{remark}

\begin{definition}[Expectation]\label{def:expectation}
Let $X$ be a random variable on $(\Omega, \cF, P)$. The \emph{expectation} (or \emph{mean}) of $X$ is the Lebesgue integral
\[
E[X] = \int_\Omega X(\omega)\, dP(\omega),
\]
provided $\int_\Omega |X|\, dP < \infty$ (i.e., $X \in L^1(\Omega, \cF, P)$). Equivalently, by the change-of-variables formula for pushforward measures,
\[
E[X] = \int_\R x\, d\mu_X(x),
\]
where $\mu_X = P \circ X^{-1}$ is the law of $X$.
\end{definition}

\begin{remark}
The expectation is a linear functional on $L^1(P)$: $E[aX + bY] = aE[X] + bE[Y]$. It is also monotone ($X \leq Y$ a.s.\ implies $E[X] \leq E[Y]$) and continuous under dominated convergence. These three properties---linearity, monotonicity, and continuity---characterize the Lebesgue integral and are the workhorses of probabilistic computation.
\end{remark}

\begin{remark}[Variance]
The \emph{variance} $\Var(X) = E\bigl[(X - E[X])^2\bigr] = E[X^2] - (E[X])^2$ measures the spread of $X$ around its mean. It is the quadratic form that appears in the CLT (\cref{thm:add-clt}): the normalization $S_n / \sqrt{n}$ is chosen precisely so that $\Var(S_n / \sqrt{n}) = \sigma^2$ for all $n$.
\end{remark}

We turn to the central construction: conditional expectation with respect to a sub-$\sigma$-algebra. This is the measure-theoretic generalization of the elementary conditional probability and requires a non-trivial existence theorem.

\begin{definition}[Conditional expectation]\label{def:cond-exp}
Let $X \in L^1(\Omega, \cF, P)$ and let $\mathcal{G} \subseteq \cF$ be a sub-$\sigma$-algebra. The \emph{conditional expectation} of $X$ given $\mathcal{G}$, denoted $E[X \mid \mathcal{G}]$, is the $P$-a.s.\ unique $\mathcal{G}$-measurable function satisfying
\begin{equation}\label{eq:cond-exp}
\int_G E[X \mid \mathcal{G}]\, dP = \int_G X\, dP \qquad \text{for all } G \in \mathcal{G}.
\end{equation}
\end{definition}

\begin{theorem}[Existence and uniqueness]\label{thm:cond-exp-exists}
For any $X \in L^1(\Omega, \cF, P)$ and any sub-$\sigma$-algebra $\mathcal{G} \subseteq \cF$, the conditional expectation $E[X \mid \mathcal{G}]$ exists and is unique $P$-a.s.
\end{theorem}

\begin{proof}
Define the signed measure $\nu(G) = \int_G X\, dP$ for $G \in \mathcal{G}$. Then $\nu \ll P|_\mathcal{G}$ (absolute continuity follows from $X \in L^1$). By the Radon--Nikod\'ym theorem, there exists a $\mathcal{G}$-measurable function $f$ with $\nu(G) = \int_G f\, dP$ for all $G \in \mathcal{G}$, unique $P$-a.s. Set $E[X \mid \mathcal{G}] = f$.
\end{proof}

\begin{remark}[Projection interpretation]
The conditional expectation $E[X \mid \mathcal{G}]$ is the orthogonal projection of $X$ onto $L^2(\Omega, \mathcal{G}, P)$ when $X \in L^2$. This geometric interpretation illuminates its fundamental properties:
\begin{enumerate}[label=(\roman*)]
\item \textbf{Tower property}: if $\mathcal{H} \subseteq \mathcal{G} \subseteq \cF$, then $E\bigl[E[X \mid \mathcal{G}] \mid \mathcal{H}\bigr] = E[X \mid \mathcal{H}]$. (Composing two projections onto nested subspaces yields the projection onto the smaller subspace.)
\item \textbf{Pulling out what is known}: if $Y$ is $\mathcal{G}$-measurable and bounded, then $E[XY \mid \mathcal{G}] = Y \cdot E[X \mid \mathcal{G}]$.
\item \textbf{Independence}: if $\sigma(X)$ and $\mathcal{G}$ are independent, then $E[X \mid \mathcal{G}] = E[X]$.
\end{enumerate}
Property~(iii) is the probabilistic content of orthogonality: independent information is uninformative.
\end{remark}

\begin{example}[La Gioconda and the lattice of $\sigma$-algebras]\label{ex:gioconda}
Let $\Omega$ be the set of all paintings entitled ``Mona Lisa'' (La Gioconda)---Leonardo's original at the Louvre, the near-contemporaneous copy at the Museo del Prado, and the countless reproductions produced over five centuries. Equip $\Omega$ with the uniform probability measure and consider three nested sub-$\sigma$-algebras:
\[
\mathcal{H} \;\subseteq\; \mathcal{G} \;\subseteq\; \cF,
\]
where:
\begin{enumerate}[label=(\roman*)]
\item $\mathcal{H} = \sigma(\mathrm{file\; format})$: the coarsest level---an observer sees only that the image is a \texttt{.jpg} file at $960$\,px resolution. At this resolution, Leonardo's original and the Prado copy are \emph{indistinguishable}: they belong to the same $\mathcal{H}$-atom.
\item $\mathcal{G} = \sigma(\mathrm{author})$: the intermediate level---the observer knows who painted each work. Conditioning on the event $\{\mathrm{painter} = \text{Leonardo}\}$ now picks out a single element of $\Omega$. \emph{The definite article ``the'' in ``THE Mona Lisa'' is a conditional operation}: it is well-defined only relative to a $\sigma$-algebra fine enough to separate the original from its copies.
\item $\cF$: full provenance---canvas material, pigment chemistry, exhibition history. Every painting is distinguished.
\end{enumerate}

The tower property reads concretely: if you first identify the painting by authorship ($\mathcal{G}$), then forget the author and retain only the file format ($\mathcal{H}$), you obtain the same answer as if you had consulted only the file format from the start. Information, once discarded, cannot be recovered by re-projection.

The moral is general: a $\sigma$-algebra is not a technicality---it is a \emph{resolution}, a choice of what distinctions the observer can make. Conditional expectation $E[X \mid \mathcal{G}]$ returns the best approximation to $X$ that is visible at the resolution~$\mathcal{G}$. The lattice of sub-$\sigma$-algebras is the lattice of all possible levels of knowledge, and the tower property asserts its internal consistency. In this light, probability theory is not primarily a theory of ``randomness''---it is a theory of \emph{information}. Randomness is not in the phenomenon; it is in the $\sigma$-algebra of the observer.
\end{example}

\begin{example}[Conditional expectation recovers conditional probability]\label{ex:cond-exp-partition}
Let $\{B_n\}$ be a countable partition of $\Omega$ into events of positive probability, and let $\mathcal{G} = \sigma(\{B_n\})$ be the $\sigma$-algebra generated by the partition. Then
\[
E[X \mid \mathcal{G}](\omega) = \sum_n \frac{E[X \mathbf{1}_{B_n}]}{P(B_n)}\, \mathbf{1}_{B_n}(\omega) = \sum_n E[X \mid B_n]\, \mathbf{1}_{B_n}(\omega).
\]
In particular, taking $X = \mathbf{1}_A$ recovers the conditional probability: $E[\mathbf{1}_A \mid \mathcal{G}] = \sum_n P(A \mid B_n)\, \mathbf{1}_{B_n}$.
\end{example}

\begin{remark}[Measure-theoretic Bayes' rule]\label{rmk:bayes-measure}
Bayes' rule (\cref{thm:bayes}) extends to the measure-theoretic setting as follows. Let $P$ and $Q$ be two probability measures on $(\Omega, \cF)$ with $Q \ll P$, and let $L = dQ/dP$ be the Radon--Nikod\'ym derivative (the \emph{likelihood ratio}). Then for any sub-$\sigma$-algebra $\mathcal{G} \subseteq \cF$ and any bounded $\cF$-measurable $X$,
\begin{equation}\label{eq:abstract-bayes}
E_Q[X \mid \mathcal{G}] = \frac{E_P[XL \mid \mathcal{G}]}{E_P[L \mid \mathcal{G}]} \qquad Q\text{-a.s.}
\end{equation}
This is the abstract Bayes formula: it relates the conditional expectation under one measure to that under another, mediated by the likelihood ratio. The elementary Bayes' rule \eqref{eq:bayes} is the special case where $\mathcal{G} = \sigma(\{B_n\})$ and $Q(A) = P(A \mid B)$ for some fixed $B$.
\end{remark}

\begin{proposition}[Optimality of conditional expectation]\label{prop:cond-exp-optimal}
Let $Y \in L^2(\Omega, \cF, P)$ and $\mathcal{G} \subseteq \cF$ a sub-$\sigma$-algebra. Among all $\mathcal{G}$-measurable predictors $f \in L^2(\Omega, \mathcal{G}, P)$, the conditional expectation $E[Y \mid \mathcal{G}]$ uniquely minimizes the mean squared error:
\begin{equation}\label{eq:mse-decomp}
E[(Y - f)^2] = E[(Y - E[Y \mid \mathcal{G}])^2] + E[(E[Y \mid \mathcal{G}] - f)^2].
\end{equation}
\end{proposition}

\begin{proof}
Write $Y - f = (Y - E[Y \mid \mathcal{G}]) + (E[Y \mid \mathcal{G}] - f)$. The first term $\varepsilon = Y - E[Y \mid \mathcal{G}]$ is orthogonal to every element of $L^2(\Omega, \mathcal{G}, P)$ by the defining property of conditional expectation: for any $\mathcal{G}$-measurable $Z \in L^2$,
\[
E[\varepsilon \, Z] = E[(Y - E[Y \mid \mathcal{G}]) \, Z] = E[YZ] - E[E[Y \mid \mathcal{G}] \, Z] = 0.
\]
The second term $E[Y \mid \mathcal{G}] - f$ is $\mathcal{G}$-measurable, hence in $L^2(\Omega, \mathcal{G}, P)$. By Pythagoras in the Hilbert space $L^2(P)$:
\[
\|Y - f\|_2^2 = \|\varepsilon\|_2^2 + \|E[Y \mid \mathcal{G}] - f\|_2^2,
\]
which is \eqref{eq:mse-decomp}. The second term vanishes if and only if $f = E[Y \mid \mathcal{G}]$ a.s., giving uniqueness.
\end{proof}

\begin{corollary}[Monotonicity]\label{cor:monotonicity}
If $\mathcal{G}_1 \subseteq \mathcal{G}_2 \subseteq \cF$, then $E[(Y - E[Y \mid \mathcal{G}_1])^2] \geq E[(Y - E[Y \mid \mathcal{G}_2])^2]$. Finer $\sigma$-algebra $\to$ lower irreducible error. This is the mathematical content of the \emph{policy improvement theorem}: finer $\sigma$-algebra $\to$ lower projection residual.
\end{corollary}

\begin{proof}
Since $\mathcal{G}_1 \subseteq \mathcal{G}_2$, we have $L^2(\Omega, \mathcal{G}_1, P) \subseteq L^2(\Omega, \mathcal{G}_2, P)$. Projecting onto a larger closed subspace can only decrease the residual.
\end{proof}

\begin{definition}[The Generalized Policy Iteration (GPI) category]\label{def:gpi-category}
Let $(\Omega, \cF, P)$ be a probability space and $Y \in L^2(\Omega, \cF, P)$.
The \emph{GPI category} $\Proj_Y$ is the thin category (poset category) whose
objects are sub-$\sigma$-algebras $\mathcal{G} \subseteq \cF$ and whose unique
morphism $\mathcal{G}_1 \to \mathcal{G}_2$ exists if and only if
$\mathcal{G}_1 \subseteq \mathcal{G}_2$.

The \emph{residual functor} is the contravariant functor
\[
R_Y : \Proj_Y^{\mathrm{op}} \to (\R_{\geq 0}, \geq), \qquad
R_Y(\mathcal{G}) = \|Y - E[Y \mid \mathcal{G}]\|_2^2.
\]
By \cref{cor:monotonicity}, $R_Y$ is a functor: every morphism
$\mathcal{G}_1 \hookrightarrow \mathcal{G}_2$ in $\Proj_Y$ maps to the
inequality $R_Y(\mathcal{G}_1) \geq R_Y(\mathcal{G}_2)$ in
$(\R_{\geq 0}, \geq)$.

The \emph{evaluation functor} is
\[
\mathcal{E}_Y : \Proj_Y \to L^2(\Omega, \cF, P), \qquad
\mathcal{E}_Y(\mathcal{G}) = E[Y \mid \mathcal{G}].
\]
A \emph{filtration} $\mathbb{F} = (\mathcal{G}_t)_{t \geq 0}$ with
$\mathcal{G}_0 \subseteq \mathcal{G}_1 \subseteq \cdots$ is a functor
$\mathbb{F} : (\mathbb{N}, \leq) \to \Proj_Y$ from the poset category of
natural numbers.
\end{definition}

\begin{example}[Generalized policy iteration and the $\sigma$-algebra lattice]\label{ex:ml-sigma}
We show that the central constructions of reinforcement learning and statistical learning are operations in the GPI category $\Proj_Y$ (\cref{def:gpi-category}). The core objects are the evaluation functor $\mathcal{E}_Y$ (\cref{def:gpi-category}) and an improvement operator $\mathcal{I}$, whose alternation is generalized policy iteration (GPI), a residual-monotone walk on $\Proj_Y$ governed by \cref{prop:cond-exp-optimal}.

\begin{enumerate}[label=(\alph*)]

\item \textbf{The two operators.} Let $Y \in L^2(\Omega, \cF, P)$ be a target random variable (a reward, a label, a future state). Let $\Lambda \subseteq \{\mathcal{G} : \mathcal{G} \subseteq \cF\}$ be the family of sub-$\sigma$-algebras reachable by the agent or model.

\emph{Evaluation operator} $\mathcal{E}$: given a fixed $\mathcal{G} \in \Lambda$, compute the orthogonal projection
\[
\mathcal{E}_{\mathcal{G}}[Y] = E[Y \mid \mathcal{G}].
\]
By \cref{prop:cond-exp-optimal}, this is the unique minimizer of $E[(Y - f)^2]$ over $f \in L^2(\Omega, \mathcal{G}, P)$. The $\sigma$-algebra is fixed; the function moves.

\emph{Improvement operator} $\mathcal{I}$: given the current projection residual, select a new $\sigma$-algebra $\mathcal{G}' \in \Lambda$ with smaller residual:
\[
\mathcal{I}(\mathcal{G}) = \mathcal{G}' \in \Lambda \quad \text{such that} \quad E[(Y - E[Y \mid \mathcal{G}'])^2] \leq E[(Y - E[Y \mid \mathcal{G}])^2].
\]
The target is fixed; the $\sigma$-algebra moves.

Note that $\mathcal{I}$ selects a $\sigma$-algebra with smaller \emph{residual}; this does not require $\mathcal{G} \subseteq \mathcal{G}'$ (the new $\sigma$-algebra need not be finer). When the walk is additionally a filtration ($\mathcal{G}_n \subseteq \mathcal{G}_{n+1}$ for all $n$), the stronger convergence of \cref{thm:in-context-gpi} applies---this is the case for in-context learning, where each new token extends the context.

\item \textbf{Generalized policy iteration.} GPI alternates the two operators:
\[
\mathcal{G}_0 \xrightarrow{\;\mathcal{E}\;} V_0 \xrightarrow{\;\mathcal{I}\;} \mathcal{G}_1 \xrightarrow{\;\mathcal{E}\;} V_1 \xrightarrow{\;\mathcal{I}\;} \mathcal{G}_2 \xrightarrow{\;\mathcal{E}\;} \cdots
\]
where $V_n = E[Y \mid \mathcal{G}_n]$. By \cref{cor:monotonicity}, each improvement step yields
\[
\|Y - V_{n+1}\|_2^2 \leq \|Y - V_n\|_2^2.
\]
The sequence of projection residuals is non-increasing and bounded below by zero, hence converges. In the finite-state setting, this convergence is achieved in finitely many steps and yields the optimal value function $V^* = E[Y \mid \mathcal{G}^*]$. This is the measure-theoretic skeleton of the \emph{policy improvement theorem}: the monotonicity of orthogonal projections on the lattice is the mechanism underlying the classical Bellman convergence argument.

\item \textbf{Reinforcement learning: the Bellman equation as tower property.} In a Markov decision process $(\mathcal{S}, \mathcal{A}, P, R, \gamma)$, the agent selects actions $A_t$ in states $S_t$, receiving rewards $R_{t+1}$. A policy $\pi$ determines the agent's behavior; the Markov property ensures
that the value function depends only on the current state, so the relevant
$\sigma$-algebra is $\sigma(S_t)$, a coarsening of the full trajectory
$\sigma$-algebra. The \emph{value function} is
\[
V^\pi(s) = E^\pi\!\left[\sum_{k=0}^\infty \gamma^k R_{t+k+1} \;\middle|\; S_t = s\right] = E[\mathcal{R}_t \mid \sigma(S_t)],
\]
where $\mathcal{R}_t = \sum_{k=0}^\infty \gamma^k R_{t+k+1}$ is the discounted return and $\gamma \in [0,1)$. This is the evaluation operator $\mathcal{E}$ applied to the return.

The \emph{Bellman equation} is the tower property in disguise:
\begin{equation}\label{eq:bellman}
V^\pi(s) = E^\pi[R_{t+1} + \gamma\, V^\pi(S_{t+1}) \mid S_t = s].
\end{equation}
Writing $\mathcal{G}_t = \sigma(S_t)$, this reads $E[\mathcal{R}_t \mid \mathcal{G}_t] = E\bigl[R_{t+1} + \gamma\, E[\mathcal{R}_{t+1} \mid \mathcal{G}_{t+1}] \mid \mathcal{G}_t\bigr]$---the tower property applied to the return, with the Markov property ensuring $\mathcal{G}_t$ captures the relevant information.

Policy improvement selects $\pi'$ greedy with respect to $V^\pi$: choose $A_t$ to maximize $E[R_{t+1} + \gamma V^\pi(S_{t+1}) \mid S_t, A_t]$. This generates a $\sigma$-algebra $\mathcal{G}_{\pi'}$ that resolves the action-value distinctions $\pi$ left unresolved, and by the same monotonicity principle, $V^{\pi'} \geq V^\pi$
(the classical policy improvement theorem; the pointwise inequality
follows from applying the Bellman operator's monotonicity state by state).

\item \textbf{Supervised learning as single-step evaluation.} The static supervised setting is the degenerate case of GPI: no dynamics, no sequential decisions. The joint distribution $(X, Y) \sim P$ is fixed. A model with parameters $\theta$ computes a representation $h_\theta : \mathcal{X} \to \R^d$ generating
\[
\mathcal{G}_\theta = \sigma(h_\theta(X)) = \{h_\theta^{-1}(B) : B \in \mathcal{B}(\R^d)\}.
\]
Two inputs $x, x'$ with $h_\theta(x) = h_\theta(x')$ are indistinguishable to the model---they belong to the same $\mathcal{G}_\theta$-atom, exactly as the Louvre and Prado Giocondas belong to the same $\mathcal{H}$-atom at 960px (\cref{ex:gioconda}). The evaluation operator gives the optimal readout $E[Y \mid \mathcal{G}_\theta]$; the improvement operator is stochastic gradient descent on $\theta$:
\begin{equation}\label{eq:ml-fundamental}
\min_\theta\, E[(Y - E[Y \mid \mathcal{G}_\theta])^2].
\end{equation}
There is no alternation---$\mathcal{E}$ and $\mathcal{I}$ are fused into a single optimization. GPI degenerates to gradient descent on the projection residual.

\item \textbf{Architecture $=$ reachable $\sigma$-algebras.} Different architectures constrain which $\sigma$-algebras are reachable:
\begin{itemize}
\item A \emph{linear model}: $\mathcal{G}_\theta = \sigma(w^\top X)$---a hyperplane partition. Coarse.
\item A \emph{CNN}: $\mathcal{G}_\theta$ respects translation equivariance---only translation-invariant distinctions are visible.
\item A \emph{Transformer} with context length $T$: $\mathcal{G}_\theta = \sigma(h_\theta(x_1, \ldots, x_T))$---the $\sigma$-algebra generated by attending to all positions in the context window.
\end{itemize}
Architecture design is the art of choosing which region of the lattice $\Lambda$ to search over.

\item \textbf{Overfitting, underfitting, and regularization.}
\begin{itemize}
\item \emph{Underfitting}: $\mathcal{G}_\theta$ is too coarse---important distinctions are invisible. The projection $E[Y \mid \mathcal{G}_\theta]$ has large residual.
\item \emph{Overfitting}: $\mathcal{G}_\theta$ is too fine on training data---it separates noise from noise. On test data from the same $P$, these spurious distinctions do not hold.
\item \emph{Regularization} (dropout, weight decay, early stopping) $=$ constraining $\Lambda$, preventing $\mathcal{G}_\theta$ from becoming too fine relative to the sample size.
\end{itemize}

\item \textbf{Next-token prediction (GPT).} The sequence $(X_1, X_2, \ldots)$ is drawn from $P$. The autoregressive model computes
\[
f_\theta(x_1, \ldots, x_t) \approx E[X_{t+1} \mid \sigma(X_1, \ldots, X_t)]
\]
for each $t$. This is GPI with a rolling evaluation: at each position, $\mathcal{E}$ projects the next token onto the $\sigma$-algebra generated by the context, and $\mathcal{I}$ updates $\theta$ to reduce the cross-entropy (the KL divergence from the model's conditional to the true conditional). The transformer's attention mechanism parameterizes the map from context to the $\sigma$-algebra $\mathcal{G}_\theta^{(t)}$ that determines what the model ``remembers.''

\item \textbf{World models (JEPA).} A world model predicts future states given current observations. In the framework:
\[
\text{JEPA predictor} = E[h_\theta(X_{t+1}) \mid \mathcal{G}_\theta^{(t)}].
\]
The key design choice: predict in \emph{representation space} $h_\theta(X_{t+1})$, not in pixel space $X_{t+1}$. This means the target itself lives in a learned $\sigma$-algebra. Both $\mathcal{E}$ and its target co-evolve---the model is simultaneously choosing what to predict and how to predict it. This is joint optimization over two $\sigma$-algebras: GPI where the improvement operator $\mathcal{I}$ moves both the predictor and the target.

\item \textbf{Diffusion models.} The forward process $X_0 \to X_1 \to \cdots \to X_T$ is a Markov chain that progressively destroys information: the mutual information $I(X_0; X_t)$ decreases with $t$, so the conditional expectation $E[f(X_0) \mid X_t]$ becomes coarser as $t$ grows. At $t = T$, the conditional distribution is nearly independent of $X_0$ (pure noise). The reverse (generative) process recovers information step by step---this is GPI run backwards on the lattice. The score function $\nabla_{x_t} \log p(x_t)$ is the gradient of the log-density, and score matching learns to invert each information-destroying step. Generation $=$ climbing from uninformative to informative, one evaluation--improvement pair at a time.

\end{enumerate}
\end{example}

\begin{remark}[The GPI--measure-theory dictionary]\label{rmk:ml-dictionary}
The following table summarizes the translation developed in \cref{ex:ml-sigma}.

\medskip
\begin{center}
\renewcommand{\arraystretch}{1.25}
\small
\begin{tabular}{@{}p{0.38\textwidth}p{0.57\textwidth}@{}}
\hline
\textbf{AI concept} & \textbf{Measure-theoretic object} \\
\hline
Target (reward, label, future state) $Y$ & Random variable in $L^2(\Omega, \cF, P)$ \\
Evaluation operator $\mathcal{E}$ & Orthogonal projection $E[Y \mid \mathcal{G}]$ \\
Improvement operator $\mathcal{I}$ & Lattice navigation $\mathcal{G} \mapsto \mathcal{G}'$ \\
GPI convergence & Monotone walk on $\Lambda$, bounded $\Rightarrow$ converges \\
Bellman equation & Tower property of conditional expectation \\
Policy improvement theorem & \cref{cor:monotonicity}: finer $\mathcal{G} \to$ smaller residual \\
Value function $V^\pi$ & $E[\mathcal{R}_t \mid \sigma(S_t)]$ (projection of return) \\
\hline
Learned representation $h_\theta(X)$ & Generator of $\mathcal{G}_\theta = \sigma(h_\theta(X))$ \\
Architecture constraints & Reachable sub-lattice $\Lambda \subseteq \{\mathcal{G} \subseteq \cF\}$ \\
Supervised training loss & $E[(Y - E[Y \mid \mathcal{G}_\theta])^2]$ (projection residual) \\
Training over $\theta$ & Search over $\sigma$-algebras in $\Lambda$ \\
Overfitting & $\mathcal{G}_\theta$ too fine (separates noise) \\
Regularization & Constraining $\Lambda$ \\
\hline
Next-token prediction & $E[X_{t+1} \mid \sigma(X_1, \ldots, X_t)]$ \\
World model (JEPA) & $E[h_\theta(X_{t+1}) \mid \mathcal{G}_\theta^{(t)}]$ (joint $\sigma$-algebra opt.) \\
Diffusion forward & Information destruction: $I(X_0; X_T) \ll I(X_0; X_0)$ \\
Diffusion reverse & Information recovery: climbing toward finer conditional structure \\
Generalization & $\mathcal{G}_\theta$ captures structure of $P$, not artifacts \\
\hline
GPI category $\Proj_Y$ & Thin category $(\Lambda, \subseteq)$ (\cref{def:gpi-category}) \\
In-context learning & Filtration $\mathbb{F} : (\mathbb{N}, \leq) \to \Proj_Y$ \\
In-context convergence & Colimit $\mathcal{G}_\infty$ in $\Proj_Y$ (\cref{thm:in-context-gpi}) \\
\hline
\end{tabular}
\normalsize
\end{center}
\medskip

\noindent Every model, every architecture, every training procedure in modern AI is a particular strategy for navigating the lattice of sub-$\sigma$-algebras of $\cF$. Reinforcement learning makes the navigation explicit through the alternation $\mathcal{E} \circ \mathcal{I}$; supervised learning fuses the two into a single gradient step. The lattice itself---and the projection theorem that governs it---is Kolmogorov's, 1933.

\smallskip
\noindent\textit{Author's note.} Machine learning, as practiced, does not resemble the popular image of ``artificial intelligence''---it resembles $\sigma$-algebra selection. The objects manipulated by gradient descent are not thoughts, beliefs, or goals; they are partitions of a sample space. The questions answered by a trained model are not ``what should I do?''\ but ``which distinctions does the data support?''\ The $\sigma$-algebra lattice is a more faithful portrait of what these systems are and do than any anthropomorphic metaphor.
\end{remark}

\begin{theorem}[In-context GPI]\label{thm:in-context-gpi}
Let $Y \in L^2(\Omega, \cF, P)$ and let
$\mathbb{F} : (\mathbb{N}, \leq) \to \Proj_Y$ be a filtration, with
$\mathbb{F}(t) = \mathcal{G}_t$ and
$\mathcal{G}_0 \subseteq \mathcal{G}_1 \subseteq \cdots \subseteq \cF$.
\begin{enumerate}[label=(\roman*)]
\item \emph{(Monotone residual.)} The composite
$R_Y \circ \mathbb{F}^{\mathrm{op}}$ is order-preserving:
$R_Y(\mathcal{G}_t)$ is non-increasing in $t$.

\item \emph{(Colimit convergence.)} Let
$\mathcal{G}_\infty = \sigma\!\bigl(\bigcup_{t \geq 0} \mathcal{G}_t\bigr)$
be the colimit of $\mathbb{F}$ in $\Proj_Y$. Then
\[
\mathcal{E}_Y(\mathcal{G}_t) = E[Y \mid \mathcal{G}_t]
\xrightarrow{\;L^2\;} E[Y \mid \mathcal{G}_\infty]
= \mathcal{E}_Y(\mathcal{G}_\infty)
\quad \text{as } t \to \infty.
\]
The evaluation functor sends the colimit of the filtration to the
$L^2$-limit of the projections.

\item \emph{(Terminal projection.)}
$R_Y(\mathcal{G}_\infty) = \inf_{t \geq 0} R_Y(\mathcal{G}_t)$.
The colimit achieves the infimum of the residual functor over the image
of $\mathbb{F}$.
\end{enumerate}
\end{theorem}

\begin{proof}
(i) is \cref{cor:monotonicity} applied to each inclusion
$\mathcal{G}_t \subseteq \mathcal{G}_{t+1}$.

(ii) The sequence $V_t = E[Y \mid \mathcal{G}_t]$ is a martingale
adapted to $(\mathcal{G}_t)$: for $s \leq t$,
$E[V_t \mid \mathcal{G}_s] = E[Y \mid \mathcal{G}_s] = V_s$ by the tower
property. Since $\sup_t E[V_t^2] \leq E[Y^2] < \infty$, L\'{e}vy's upward
theorem \cite[Theorem~4.4.5]{Durrett} gives
$V_t \to E[Y \mid \mathcal{G}_\infty]$ both a.s.\ and in $L^2$.

(iii) By (ii) and continuity of the $L^2$-norm:
$R_Y(\mathcal{G}_t) = \|Y - V_t\|_2^2 \to
\|Y - E[Y \mid \mathcal{G}_\infty]\|_2^2 = R_Y(\mathcal{G}_\infty)$.
Combined with (i), $R_Y(\mathcal{G}_\infty)$ is the infimum.
\end{proof}

\begin{remark}[In-context reinforcement learning]\label{rmk:in-context-rl}
\cref{thm:in-context-gpi} gives the measure-theoretic content of
\emph{in-context reinforcement learning}: a transformer processing a
sequence of $(s_k, a_k, r_{k+1})$ tuples builds a filtration
$\mathcal{G}_t = \sigma(s_1, a_1, r_2, \ldots, s_t)$ within its context
window. Each new token is a morphism
$\mathcal{G}_t \hookrightarrow \mathcal{G}_{t+1}$ in $\Proj_Y$. The three
modes of GPI are distinguished by the role of the improvement operator:

\medskip
\begin{center}
\renewcommand{\arraystretch}{1.25}
\small
\begin{tabular}{@{}p{0.22\textwidth}p{0.34\textwidth}p{0.36\textwidth}@{}}
\hline
\textbf{Mode} & \textbf{Improvement $\mathcal{I}$} & \textbf{Categorical structure} \\
\hline
Classical RL & Explicit: $\pi \to \pi'$ & Endofunctor on $\Proj_Y$ \\
Supervised learning & Fused: $\nabla_\theta$ on $R_Y(\mathcal{G}_\theta)$ & Search over $\Lambda \subset \mathrm{Ob}(\Proj_Y)$ \\
In-context RL & Implicit: data extends context & Filtration $\mathbb{F} : (\mathbb{N}, \leq) \to \Proj_Y$ \\
\hline
\end{tabular}
\normalsize
\end{center}
\medskip

\noindent In the in-context case, the improvement operator is the identity
on the filtration---improvement is free, supplied by the data. The
transformer's forward pass at position $t$ approximates the evaluation
functor $\mathcal{E}_Y(\mathcal{G}_t)$, and convergence is guaranteed by
\cref{thm:in-context-gpi}(ii).
\end{remark}

\begin{remark}[Connection to the CLT framework]
The foundations reviewed here underlie every result in this paper. The CLT (\cref{thm:add-clt}) concerns the law $\mu_{S_n/\sqrt{n}} = P \circ (S_n/\sqrt{n})^{-1}$ of a normalized sum, whose characteristic function is an expectation: $\varphi_{S_n/\sqrt{n}}(\xi) = E[e^{i\xi S_n/\sqrt{n}}]$. The Fourier transform of a probability measure is nothing but the expectation of a character. The entire duality between the additive and multiplicative CLTs (\cref{thm:mellin-fourier}) rests on the measure-theoretic pushforward: if $\mu$ is a measure on $(\Rpos, \times)$, then $\log_*\mu$ is a measure on $(\R, +)$, and $\cM[\mu](s) = \cF[\log_*\mu](s)$ (\cref{thm:mellin-fourier}). The logarithm is a group isomorphism, not a linear map; the key property is that it preserves Haar measure ($dx/x \mapsto dy$), not that it commutes with expectation.
\end{remark}

